% ============================================================
% question1.tex - 问题一的模型建立与求解
% 每一问使用相同结构，复制此文件为 question2.tex 等
% ============================================================

\section{问题一的模型建立与求解}

% ==============================
\subsection{方法对比与选择}
% ==============================

% 列出至少 2 种可行方法，用对比表展示优劣，
% 正文中对每种方法简要描述原理（各 100-200 字），
% 最后明确说明选择理由。

针对问题一，本文考虑了以下候选方法：

\begin{table}[H]
  \centering
  \caption{问题一候选方法对比}
  \begin{tabular}{lcccc}
    \toprule
    \textbf{方法} & \textbf{适用性} & \textbf{精度} & \textbf{效率} & \textbf{可解释性} \\
    \midrule
    方法A & 高 & 高 & 中 & 高 \\
    方法B & 中 & 高 & 高 & 低 \\
    \bottomrule
  \end{tabular}
\end{table}

（对每种方法的简要描述...）

综合考虑适用性和可解释性，本文选择方法A求解问题一。

% ==============================
\subsection{模型建立}
% ==============================

% 目标函数、约束条件、完整数学表达。

\subsubsection{决策变量}

（变量定义）

\subsubsection{目标函数}

\begin{equation}
  \min Z = \sum_{i=1}^{n} c_i x_i
  \label{eq:q1_obj}
\end{equation}

\subsubsection{约束条件}

\begin{equation}
  \text{s.t.} \quad
  \begin{cases}
    \sum_{i=1}^{n} a_{ij} x_i \leq b_j, & j = 1, 2, \ldots, m \\
    x_i \geq 0, & i = 1, 2, \ldots, n
  \end{cases}
  \label{eq:q1_cons}
\end{equation}

% ==============================
\subsection{求解过程}
% ==============================

% 算法步骤、参数设置、收敛条件。

（求解算法描述、参数设置说明）

% ==============================
\subsection{结果与分析}
% ==============================

% 表格：原始计算数据 / 结果对比
\begin{table}[H]
  \centering
  \caption{问题一求解结果}
  \begin{tabular}{ccc}
    \toprule
    \textbf{指标} & \textbf{数值} & \textbf{单位} \\
    \midrule
    目标函数值 & 1256.3 & 万元 \\
    求解时间   & 2.4    & 秒 \\
    \bottomrule
  \end{tabular}
\end{table}

% 图：至少 4 张，每张至少 100 字分析

\begin{figure}[H]
  \centering
  \includegraphics[width=0.8\textwidth]{q1_fig1.pdf}
  \caption{问题一结果分析图1}
  \label{fig:q1_1}
\end{figure}

如图~\ref{fig:q1_1}所示，（至少100字的详细分析...）

\begin{figure}[H]
  \centering
  \includegraphics[width=0.8\textwidth]{q1_fig2.pdf}
  \caption{问题一结果分析图2}
  \label{fig:q1_2}
\end{figure}

如图~\ref{fig:q1_2}所示，（至少100字的详细分析...）

% 按需添加更多图（每问至少 4 张）

% ---- 灵敏度/稳定性分析 ----
\subsubsection{灵敏度分析}

（对关键参数的灵敏度分析结果）
